% Required packages in the main preamble:
% \usepackage{graphicx}
% \usepackage{booktabs}
% \usepackage{tabularx}
% \usepackage{array}

\providecommand{\SysName}{CARE}

\section{System Design}
\label{sec:system-design}

\SysName{} is designed around a simple principle: LLM output is useful for
proposing source edits, but it must not be trusted as evidence of correctness.
The system therefore separates program understanding, opportunity detection,
LLM-based patch generation, and validation into explicit stages. Each stage
produces machine-readable artifacts that constrain the next stage and provide
traceability from an accepted patch back to the evidence that motivated it.

\subsection{Architecture and Design Goals}
\label{sec:design-architecture}

Figure~\ref{fig:system-design-overview} shows the high-level design. A target
C/C++ project enters through a project loader that discovers source files,
build-system metadata, optional compile databases, and user-supplied build and
test commands. \SysName{} then constructs a security-aware context graph, runs
opportunity detectors over that graph, converts selected opportunities into
explicit safety contracts, asks an LLM to propose patch candidates, and accepts a
candidate only after the validation pipeline succeeds.

\begin{figure}[t]
  \centering
  \includegraphics[width=\columnwidth]{img/care_system_design_overview.pdf}
  \caption{Single-column overview of \SysName{}. Program context and security
  knowledge constrain detection and LLM patch generation. Validation failures
  are summarized as feedback for smaller repair prompts, and a patch is selected
  only after all required gates pass.}
  \label{fig:system-design-overview}
\end{figure}

Three goals drive the design. First, \SysName{} must preserve security-sensitive
behavior, including cleanup ordering, return-value conventions, validation
checks, and externally visible side effects. Second, it must remain usable on
real C/C++ projects whose compile databases, generated headers, or build systems
may be incomplete. Third, it must make each result auditable: every patch should
be linked to an opportunity, detector evidence, a refactoring plan, validation
logs, and failure explanations for rejected candidates.

The architecture is intentionally fail-closed. A stage may increase confidence
in a candidate, but no stage is allowed to accept a patch by itself. For
example, an LLM-generated diff that matches the planner's intent is still
discarded if it does not apply cleanly, fails the project build, changes a
public signature unexpectedly, removes a guard, or creates an obvious resource
imbalance. This design makes LLM generation a proposal mechanism rather than a
correctness oracle.

\begin{table*}[t]
  \centering
  \caption{Stage contracts in \SysName{}. Each stage emits structured artifacts
  consumed by later stages and used for reporting.}
  \label{tab:system-stage-contracts}
  \footnotesize
  \begin{tabularx}{\textwidth}{@{}>{\raggedright\arraybackslash}p{0.16\textwidth}>{\raggedright\arraybackslash}p{0.23\textwidth}>{\raggedright\arraybackslash}p{0.25\textwidth}>{\raggedright\arraybackslash}X@{}}
    \toprule
    \textbf{Stage} & \textbf{Primary input} & \textbf{Structured output} & \textbf{Safety role} \\
    \midrule
    Project loading &
    Source tree, build/test commands, compile database &
    Source-file list, build metadata, clean backup/restore handle &
    Keeps patch validation isolated and repeatable. \\
    Context construction &
    Parsed functions and project metadata &
    Context graph with functions, blocks, typed edges, resource events, summaries &
    Gives detectors and prompts path-, resource-, call-, and side-effect context. \\
    Detection and planning &
    Context graph plus security knowledge &
    Ranked opportunities and JSON-safe refactoring plans &
    Converts a vague cleanup idea into explicit invariants and constraints. \\
    Patch generation &
    Plan, source snippet, full function body, optional feedback &
    One or more unified-diff candidates or JSON edit plans converted to diffs &
    Limits the LLM to minimal edits scoped to affected files and functions. \\
    Validation and repair &
    Candidate diff applied to a clean tree &
    Per-stage validation result, failure taxonomy, selected patch or repair feedback &
    Rejects unsafe candidates and feeds concrete failures back to the generator. \\
    Reporting &
    Opportunities, plans, patches, validation logs &
    JSON/Markdown reports and per-candidate diffs &
    Preserves traceability from finding to accepted or rejected patch. \\
    \bottomrule
  \end{tabularx}
\end{table*}

Table~\ref{tab:system-stage-contracts} summarizes the contract between stages.
The important design decision is that intermediate artifacts are explicit rather
than hidden inside prompts. A detector does not merely say that code ``looks
duplicated''; it emits source locations, duplicated statements, resource
variables, affected paths, severity, confidence, and a context summary. Likewise,
the planner does not directly produce a patch. It produces an intent, affected
files, required invariants, security constraints, behavior-preservation goals,
and a validation strategy.

This explicit artifact flow also supports incremental experimentation. The same
detected opportunity can be sent to different LLMs, generated in conservative
or aggressive mode, validated with different ablations, or inspected by a human
reviewer without re-running the entire front end. In our prototype, all major
artifacts are serializable to JSON so that experiments can be resumed, audited,
and aggregated across projects.

\subsection{Context Construction}
\label{sec:design-context}

The context construction stage builds the program facts needed to reason about a
candidate refactoring. \SysName{} first parses C/C++ source files into
function-level records containing signatures, bodies, line ranges, local
variables, returns, labels, gotos, and calls. When clang and compile arguments
are available, the parser uses clang's JSON AST to extract real translation-unit
facts; otherwise it falls back to a deterministic brace-aware parser. This
fallback is intentional: large experiments often include projects that cannot be
fully compiled in a fresh environment, but their source can still provide useful
local refactoring context.

For each function, \SysName{} constructs a lightweight local CFG,
extracts call edges, and computes local data-flow facts. The CFG
captures sequential blocks, branches, loops, returns, gotos, and labels. The
data-flow pass records definitions, uses, pointer variables, return-status
variables, and resource-like variables. A security knowledge base then matches
resource events such as allocation/free, open/close, and lock/unlock pairs, as
well as security smells such as unchecked return values, unsafe APIs, repeated
checks, and inconsistent error propagation.

These facts are merged into a security-aware contextual refactoring graph. The
graph contains function and basic-block nodes; call, control, data, goto,
return, and resource edges; resource events; error-handling path summaries;
side-effect summaries; and return-condition summaries. \SysName{} uses this graph
in two ways. Detectors use the structured graph to avoid purely textual matches,
and prompt builders use concise graph summaries to keep LLM context small but
safety-relevant. For example, a cleanup opportunity prompt can include which
paths return early, which variables are resources, which calls have visible side
effects, and which release operations already exist.

The graph has two levels of granularity. Function-level records provide stable
units for ranking, prompting, and reporting, while basic-block records preserve
local control-flow structure needed by cleanup and reachability detectors. Edges
are typed rather than merely labeled strings: call edges connect caller blocks
to known callees, control edges encode intra-procedural reachability, data edges
summarize definition-use relationships, goto edges connect jumps to labels, and
resource edges connect acquire and release events. This schema is lightweight,
but it gives later stages enough structure to distinguish a redundant-looking
check from one that follows a possible mutation or error-handling side effect.

The graph is deliberately conservative. It does not attempt to prove complete
C/C++ semantics or deep aliasing properties. Instead, it records the facts that
are most useful for safe local refactoring and exposes uncertainty through
evidence confidence, post-filter tags, and optional clang confirmation. This
lets \SysName{} rank opportunities without pretending that every source-level
pattern is equally safe to edit.

Context summaries are generated from the same graph rather than written by the
detectors ad hoc. A function summary includes the signature, return conventions,
resource lifecycle, visible calls, branch and cleanup structure, and detector
evidence. A location summary narrows this information to the surrounding block,
incoming and outgoing control paths, and resource state near the target line.
These summaries are short by design: they give the LLM the facts that matter for
the proposed transformation while keeping irrelevant project text out of the
prompt.

\subsection{Detection and Planning}
\label{sec:design-detection-planning}

Opportunity detection is organized as a set of function-level detectors sharing
the same interface: each detector receives a function, the context graph, and
the security knowledge base, then returns zero or more refactoring
opportunities. \SysName{} currently includes detectors for cleanup-oriented goto
patterns, resource imbalance, duplicate checks, dead or unreachable code, and
general security smells. The detector orchestrator deduplicates findings,
clusters related resource findings, applies post-filters, optionally asks clang
to confirm high-impact findings, and ranks results by severity and confidence.

Ranking combines impact and evidence quality. Impact is higher when the finding
touches resource lifecycle, error propagation, externally visible side effects,
or security checks. Evidence quality is higher when source locations are precise,
the relevant resource or checked variable is identified, the affected paths are
known, and no conflicting mutation or aliasing signal is observed. This prevents
cosmetic refactorings from crowding out high-value safety fixes, while still
allowing low-risk cleanup opportunities to appear later in the queue.

The detectors emit evidence rather than patches. A resource detector, for
instance, records the resource variable, acquire and
release events, missing or repeated release operations, early-return locations,
and the expected matching cleanup operation. A duplicate-check detector records
the checked variable, the first and later checks, and whether an assignment,
mutating call, or obvious aliasing risk appears between them. This evidence is
used both for ranking and for constraining the LLM.

Planning turns each selected opportunity into a safety contract. The planner
derives affected files and functions, required invariants, security constraints,
behavior-preservation goals, a patch strategy, and a validation strategy. This
step is separate from patch generation so that the system can inspect, serialize,
and report the intended transformation before asking an LLM to edit code. The
plan is also where \SysName{} encodes refactoring-specific rules: duplicate
checks may be removed only when the detector reports no intervening mutation or
aliasing risk; cleanup consolidation must preserve cleanup order and return
codes; resource repairs must avoid leaks and double releases.

The safety contract is the main boundary between program analysis and the LLM.
It does not ask the model to rediscover every invariant from raw source text.
Instead, it tells the model which invariants are non-negotiable: keep return
codes stable, preserve logging and error-reporting calls, do not change API
contracts, preserve or improve resource pairing, and do not remove validation
checks unless the detector evidence establishes redundancy. The model is then
used for what it is good at: synthesizing a small source-level rewrite that
satisfies those constraints.

The patch generator consumes the plan, source snippet, full function body,
context summary, and any validation feedback. It supports conservative and
aggressive modes, but both modes are constrained to produce small, scoped edits.
The preferred output is a JSON edit plan that \SysName{} converts into a
unified diff; a raw unified diff is accepted as a fallback. This design reduces
the chance of malformed patch paths and makes it easier to repair common LLM
formatting mistakes before validation.

Conservative mode is used for the main evaluation because it is easier to
validate: it prefers local rewrites, avoids helper extraction, and keeps edits
inside the affected function when possible. Aggressive mode may introduce helper
functions or larger cleanup abstractions when duplication spans multiple blocks
or functions. Both modes pass through the same validator, so a larger rewrite
must earn acceptance through the same build, test, semantic, resource, and
security gates as a small rewrite.

\subsection{Patch Selection and Repair}
\label{sec:design-validation}

\SysName{} treats every patch candidate as untrusted. Before validation, the
project loader creates a clean restore point using git when available, or a
temporary directory copy otherwise. The validator restores the clean tree,
repairs common patch-path mistakes, checks whether the diff applies, applies it,
and rejects candidates that produce no actual change. It then runs the build
command, the test command when provided, optional static analyzers, and three
lightweight regression checkers.

The validation order is chosen to fail fast while preserving useful diagnostics.
Patch-application errors are reported before build failures; build failures are
reported before test failures; static and semantic checks run only on an applied
candidate; and security/resource regression checks are reported with source
locations whenever possible. This order keeps expensive checks from running on
obviously invalid candidates while still collecting enough feedback for repair.

The resource checker searches for obvious missing releases, double releases,
missing unlocks, and early returns that bypass cleanup. The semantic checker
compares public function signatures, return forms, and externally visible calls
before and after the patch. The security checker rejects patches that remove
null or bounds checks unsafely, replace safer APIs with unsafe ones, introduce
unchecked allocation or return values, or remove cleanup on error paths. These
checkers are conservative: they are designed to catch common unsafe regressions,
not to prove full equivalence.

If a candidate fails, the repair loop summarizes validation logs,
counterexamples, and failure taxonomy into feedback for the next LLM call. The
repair prompt asks for a smaller, safer patch that preserves all previously
stated invariants and avoids unrelated edits. \SysName{} iterates this process
up to a configured limit and selects the first candidate that passes all
required validation stages. Failed candidates are retained in the report rather
than discarded, because their failure reasons explain what risks the validator
guarded against.

The final report contains detected opportunities, selected patches, failed
candidates, before/after summaries, validation results, security impact
summaries, performance estimates, and traceability links from each opportunity
to its patch candidates. This reporting layer is part of the design rather than
an afterthought: for security-sensitive refactoring, users need to inspect not
only what changed, but also why the change was proposed and why the system
believes it is safe enough to accept.

In summary, the system design separates evidence collection, decision making,
generation, and validation. The analysis stages provide structured context; the
detectors identify and rank opportunities; the planner turns an opportunity into
a safety contract; the LLM proposes edits under that contract; and the validator
decides whether any candidate is acceptable. This separation is what makes
\SysName{} suitable for empirical study: each stage can be measured, ablated,
or replaced without changing the rest of the pipeline.
